% supplement_submission.tex — CANDIDATE trimmed supplement, built for the author's sign-off.
%
% Not part of the submission package by default. manuscript/final/supplement_audit.md explains why:
% the author's 2026-08-25 decision was to eliminate supplementary material entirely, and this file
% proposes reversing that decision for exactly three blocks judged reviewer-useful (see the audit),
% which is a recommendation, not something applied unilaterally.
%
% Content: extracted verbatim from manuscript/drafts/supplement.tex (untouched, still in the repo as
% the analysis record) — the four Benjamini-Hochberg correction family tables, the exploratory
% proportional-odds model, and the full psychometric reliability table. Everything else in the
% original supplement (instrument wording, facilitator-practice tables, skip-logic prose, missingness
% table, qualitative scaffold, reproducibility statement) is fully superseded by v7_draft.tex content
% or was dropped by the author for a substantive (non-space) reason — see supplement_audit.md.
%
% If the author approves shipping this: relabel to S1-S6, cite it from v7's Methods §2.6 (Statistical
% approach) and Limitations, and submit alongside the main PDF as JPD's supplementary-material item
% (JPD publishes supplementary material via Figshare; text tables are not explicitly excluded by the
% live instructions page). None of that citation work is done here since the file is not yet approved.

\documentclass[12pt]{article}

\usepackage[utf8]{inputenc}
\usepackage[T1]{fontenc}
\usepackage{times}
\usepackage[margin=1in]{geometry}
\usepackage{graphicx}
\usepackage{booktabs}
\usepackage{longtable}
\usepackage{amsmath}
\usepackage{array}
\emergencystretch=3em

\renewcommand{\thetable}{S\arabic{table}}
\renewcommand{\thesection}{S\arabic{section}}

\title{Supplementary Material\\
\large Ceremonial Entheogen Use in Puerto Rico: A Descriptive Survey of Participants and a Small
Facilitator Subsample}
\date{}

\begin{document}
\maketitle

\noindent
All results below derive from the same analytic sample
described in the main text ($N=73$; participant branch $n=67$ routed; facilitator branch $n=8$, 7
dual-role). Structural-skip and ``prefer not to answer'' conventions follow the main text's Methods
and are applied uniformly. Counts are exact throughout, including cells below five. The bias-corrected
$\varepsilon^2$ formula can return small negative values when an effect is indistinguishable from
zero; these are displayed truncated to 0.00 and marked $\dagger$, applied uniformly.

\section{Multiplicity correction: all four test families}

The preparation-to-satisfaction comparison, designated in the analysis plan as the primary comparison
before any test was run, is reported uncorrected in the main text and is excluded from every family below.
The study was not preregistered. All other bivariate tests were assigned in advance to one of four
thematically organized families and corrected jointly within family by the Benjamini--Hochberg
false-discovery-rate procedure. Two variables, non-consensual contact and screening received, recur as
predictor or outcome across three of the four families (see the safety-domain and legal-knowledge families
below), so family-wise correction does not bound the error rate for conclusions about either variable across
every test in which it appears; no q-value in any family is below 0.05, so this does not change any result.
\textbf{No test in any family was significant at $q<0.05$.}

\begin{table}[htbp]
\centering
\caption{Benjamini--Hochberg correction, bivariate family. $\dagger$ marks an $\varepsilon^2$ whose
bias-corrected value was negative and is displayed truncated to 0.00. $\ddagger$: a tie-aware permutation
check (10,000 label permutations, seed 20260715) gave $p=0.010$ for this test, the only one in this family
reaching uncorrected significance; its ``no''-background group contributes only 3 non-missing values.}
\begin{tabular}{lrrr}
\toprule
Test & raw $p$ & BH $q$ & $\varepsilon^2$ \\
\midrule
Preparation (3-level) $\rightarrow$ satisfaction (KW) & 0.120 & 0.300 & 0.04 \\
Comfort/safety by prior knowledge of facilitator background (KW) & 0.326 & 0.407 & 0.00 \\
Comfort/safety by screening received (KW) & 0.301 & 0.407 & 0.01 \\
Trust in facilitator by prior knowledge of background (KW) & $0.017^{\ddagger}$ & 0.085 & 0.10 \\
Trust in facilitator by screening received (KW) & 0.679 & 0.679 & $0.00^{\dagger}$ \\
\bottomrule
\end{tabular}
\end{table}

\begin{table}[htbp]
\centering
\caption{Benjamini--Hochberg correction, substance-stratified family.}
\begin{tabular}{lrrr}
\toprule
Test & raw $p$ & BH $q$ & effect size \\
\midrule
Satisfaction by substance group (KW) & 0.083 & 0.166 & $\varepsilon^2=0.05$ \\
Non-consensual contact by substance group (permutation $\chi^2$) & 0.486 & 0.486 & $V=0.17$ \\
\bottomrule
\end{tabular}
\end{table}

\begin{table}[htbp]
\centering
\caption{Benjamini--Hochberg correction, safety-domain family. The non-consensual contact outcome is the
lifetime item, while the three participant-side exposures crossed against it all ask about the most recent
ceremony; 3 of the 5 respondents reporting the event ever placed it at an earlier ceremony, so for those
three the exposure and the outcome describe different events. This referent mismatch applies to every row
in the upper block and is the main reason these tests are hypothesis-generating only.}
\begin{tabular}{p{8.2cm}rrl}
\toprule
Test & raw $p$ & BH $q$ & effect size \\
\midrule
Non-consensual contact $\times$ prior knowledge of facilitator background & 0.047 & 0.188 & $V=0.38$ \\
Non-consensual contact $\times$ screening received & 0.094 & 0.188 & $V=0.32$ \\
\addlinespace
Witnessed crisis signs $\times$ distress-response protocol (facilitator branch) & 0.143 & 0.191 & OR undefined \\
Witnessed crisis signs $\times$ emergency plan (facilitator branch) & 0.429 & 0.429 & OR undefined \\
\bottomrule
\end{tabular}
\end{table}

\begin{table}[htbp]
\centering
\caption{Benjamini--Hochberg correction, legal-knowledge family. Predictors are three participant-side
legal items; outcomes are disclosure to a health professional, screening received, and non-consensual
contact. All tests Kruskal--Wallis. $\dagger$ marks a truncated $\varepsilon^2$, as above.}
\begin{tabular}{lrrr}
\toprule
Test & raw $p$ & BH $q$ & $\varepsilon^2$ \\
\midrule
General legal knowledge by disclosure to a health professional & 0.258 & 0.581 & 0.01 \\
Most-recent-ceremony legal knowledge by disclosure & 0.164 & 0.495 & 0.03 \\
Most-recent-ceremony legal worry by disclosure & 0.031 & 0.283 & 0.09 \\
General legal knowledge by screening received & 0.412 & 0.741 & $0.00^{\dagger}$ \\
Most-recent-ceremony legal knowledge by screening received & 0.659 & 0.741 & $0.00^{\dagger}$ \\
Most-recent-ceremony legal worry by screening received & 0.914 & 0.914 & $0.00^{\dagger}$ \\
General legal knowledge by non-consensual contact & 0.617 & 0.741 & $0.00^{\dagger}$ \\
Most-recent-ceremony legal knowledge by non-consensual contact & 0.165 & 0.495 & 0.03 \\
Most-recent-ceremony legal worry by non-consensual contact & 0.611 & 0.741 & $0.00^{\dagger}$ \\
\bottomrule
\end{tabular}
\end{table}

\noindent
The only legal-domain comparison reaching uncorrected significance was most-recent-ceremony legal worry by
disclosure to a health professional ($H=6.92$, $p=0.031$, $\varepsilon^2=0.09$): median 1.5 (IQR 0--2) among
the 8 reporting disclosure against median 0 (IQR 0--0) among the 47 reporting none. It did not survive
correction ($q=0.283$). Two further legal items were not testable: a general legal-worry item reached only
six respondents because of a skip-logic error described in the main text's Limitations, and a facilitator
legal-knowledge item sits on the unlinked facilitator branch with no participant-side counterpart.

\section{Exploratory proportional-odds model}

This model adjusts preparation for age and facilitator trust simultaneously. The full covariate set
specified in the original protocol was not estimable: gender cells were too thin for a four-level
categorical predictor, and screening overlapped substantively with the trust and background measures
already in the model. A prespecified three-predictor fallback was used instead. The outcome was collapsed
to three ordered levels for this model only, because the raw seven-level scale has cells as small as one to
three respondents at the low end and risks quasi-complete separation at this sample size; every other table
in the main text and this supplement retains the full 0--6 scale.

\begin{table}[htbp]
\centering
\caption{Proportional-odds model for reported satisfaction. Complete-case $n=62$ of the 63 participants
with a non-missing outcome. Collapsed outcome: low/mixed (0--3) $n=9$; mostly positive (4--5) $n=19$; very
positive (6) $n=34$.}
\begin{tabular}{lrrr}
\toprule
Predictor & OR & 95\% CI & $p$ \\
\midrule
Age (years) & 0.99 & 0.94--1.03 & 0.527 \\
Any preparation vs.\ none & 5.21 & 0.84--32.41 & 0.077 \\
Trust in facilitator (0--4) & 2.20 & 1.06--4.53 & 0.033 \\
\bottomrule
\end{tabular}
\end{table}

\noindent
Omnibus likelihood-ratio test against an intercept-and-cutpoint-only model: $\chi^2(3)=8.68$, $p=0.034$;
McFadden pseudo-$R^2=0.07$. \textbf{Power caveat.} Sixty-two complete cases support five estimated
parameters (three slopes and two cutpoints), roughly 12 cases per parameter, below conventional rules of
thumb for stable ordinal-regression estimates. The proportional-odds assumption was not formally tested,
the sample being too small to test it reliably, and is simply assumed. Every estimate in this table is
exploratory and hypothesis-generating, and should not be presented without this caveat.

\section{Psychometric evidence for the instrument}

The instrument was purpose-built, so this is its first reliability evidence. Two genuine same-respondent pairs are testable.

\begin{table}[htbp]
\centering
\caption{Reliability of the two testable same-respondent item pairs.}
\begin{tabular}{lrrr}
\toprule
Scale & Cronbach's $\alpha$ & 95\% CI & $n$ \\
\midrule
Trust in facilitator + in-ceremony comfort/safety (2 items) & 0.58 & 0.31--0.74 & 64 \\
Legal knowledge, general + most-recent ceremony (2 items) & 0.86 & 0.75--0.92 & 55 \\
\bottomrule
\end{tabular}
\end{table}

\noindent
With only two items, $\alpha$ is a monotonic transform of the pairwise correlation and the item-rest
correlations are not independent numbers. The corresponding Spearman correlations are $\rho=0.39$
($p=0.002$, $n=64$) for the trust-and-comfort pair and $\rho=0.79$ ($p<0.001$, $n=55$) for the
legal-knowledge pair; the latter two items required harmonization onto a shared 0--4 code because the
most-recent-ceremony version omits the bottom category its general counterpart carries.

A third candidate pair was untestable: a perceived-life-impact item had zero variance, with all 67
respondents answering selecting the same option, so it cannot correlate with anything and $\alpha$ is
undefined. This is reported as a data-quality finding, not a scale result.

\end{document}
